\section{  \LR{DH table}}
\label{sec:dh_table}

بناءً على التوزيع الهندسي لجمل الإحداثيات الموضحة في القسم السابق، يتم استخراج المعاملات الأربعة الخاصة بخوارزمية (\LR{DH}) لربط كل جملتين متتاليتين ($i-1$) و ($i$). تُعرّف هذه المعاملات الأربعة كما يلي:
\begin{itemize}
    \item $a_i$: الإزاحة الخطية على طول المحور $X_i$ من المحور $Z_{i-1}$ إلى المحور $Z_i$.
    \item $\alpha_i$: الإزاحة الزاوية (الفتل) حول المحور $X_i$ من المحور $Z_{i-1}$ إلى المحور $Z_i$.
    \item $d_i$: الإزاحة الخطية على طول المحور $Z_{i-1}$ من المحور $X_{i-1}$ إلى المحور $X_i$.
    \item $\theta_i$: الإزاحة الزاوية حول المحور $Z_{i-1}$ من المحور $X_{i-1}$ إلى المحور $X_i$.
\end{itemize}

بتطبيق هذه القواعد على البنية الميكانيكية للروبوت المدروس (\LR{SCARA})، تم استنتاج الوسطاء الحركية للمفاصل الثلاثة كما هو موضح في الجدول (\ref{tab:dh_parameters}). 

\begin{table}[H]
    \centering
    %\caption{DH table}
    \label{tab:dh_parameters}
    \renewcommand{\arraystretch}{1.5}
    \begin{tabular}{|c|c|c|c|c|c|}
        \hline
        \textbf{المفصل ($i$)} & \textbf{نوع المفصل} & \textbf{$\theta_i$} & \textbf{$d_i$} & \textbf{$a_i$} & \textbf{$\alpha_i$} \\
        \hline
        1 & دوراني (\LR{Revolute}) & $\theta_1$ & $0$ & $L1$ & $0$ \\
        \hline
        2 & دوراني (\LR{Revolute}) & $\theta_2$ & $0$ & $L_2$ & $180^\circ$ \\
        \hline
        3 & انسحابي (\LR{Prismatic}) & $0$ & $d_3$ & $0$ & $0$ \\
        \hline
    \end{tabular}
  
 
\end{table}

